\documentclass[12pt,a4paper]{article}
\usepackage[margin=1in]{geometry}
\usepackage[T1]{fontenc}
\usepackage{setspace}
\usepackage{amsmath,amssymb}
\usepackage{booktabs}
\usepackage{longtable}
\usepackage{xcolor}
\usepackage{hyperref}
\usepackage{fancyhdr}

\setstretch{1.18}
\setlength{\parskip}{0.45em}
\setlength{\parindent}{0pt}
\setlength{\headheight}{15pt}
\hypersetup{colorlinks=true,linkcolor=blue,citecolor=blue,urlcolor=blue}

\pagestyle{fancy}
\fancyhf{}
\lhead{CST428 - Blockchain Technologies}
\rhead{Module 3 Master Notes}
\cfoot{\thepage}

\newcommand{\examtip}[1]{%
\vspace{0.35em}
\noindent\fcolorbox{black}{gray!12}{\parbox{0.965\linewidth}{\textbf{Exam Tip:} #1}}%
\vspace{0.35em}
}

\begin{document}

\section{Module 3 Topic Map From Previous Questions}

From May 2024, June 2023, and October 2023 papers, Module 3 repeatedly asks these topics:

\begin{itemize}
\item Why consensus is needed in blockchain.
\item Crash Fault Tolerance (CFT) and Byzantine Fault Tolerance (BFT).
\item Paxos protocol and its phases.
\item Proof of Work (PoW) and Proof of Stake (PoS).
\item Bitcoin mining process and miner responsibilities.
\item Bitcoin transaction flow and transaction validation.
\item UTXO model and Bitcoin transaction structure.
\item Wallet concept and wallet types.
\item Byzantine Generals Problem, Sybil attack, DoS, and double-spend discussion.
\end{itemize}

\section{Core Theory (Explained Clearly)}

\subsection{Start Here: Why Consensus Exists}
In blockchain, there is no single bank server to say ``this transaction is final.''
So all nodes in the network need a common rule to agree on one valid ledger.
That common rule is called \textbf{consensus}.

If consensus is weak, these problems happen:
\begin{itemize}
\item one coin can be spent twice,
\item different users see different balances,
\item fake transaction history may spread,
\item trust in the system collapses.
\end{itemize}

\examtip{If a long question asks ``need for consensus'', begin with one sentence: ``Consensus keeps all distributed copies of the ledger in the same state without a central authority.''}

\subsection{Fault Types in Simple Words}

\begin{itemize}
\item \textbf{Crash fault:} node stops, hangs, or becomes silent.
\item \textbf{Byzantine fault:} node acts maliciously or sends wrong/conflicting data.
\end{itemize}

In one line: \textbf{CFT handles accidental crashes, BFT handles adversarial behavior.}

\subsection{Byzantine Generals Problem and PBFT Formula}
Byzantine Generals Problem asks: how can honest participants agree on one decision when some participants may lie?

PBFT requirement is:
\[
n \ge 3f + 1
\]
where:
\begin{itemize}
\item $f$ = maximum Byzantine nodes tolerated,
\item $n$ = total nodes needed.
\end{itemize}

Example in exam:
If Byzantine nodes are 6, then
\[
n \ge 3(6)+1=19
\]
So minimum required nodes are \textbf{19}.

\subsection{Paxos in an Easy Way}
Paxos is a \textbf{CFT} consensus protocol. It works well when failures are crashes, not malicious lies.

Roles:
\begin{itemize}
\item \textbf{Proposer:} suggests a value.
\item \textbf{Acceptor:} votes on proposed value.
\item \textbf{Learner:} learns final agreed value.
\end{itemize}

Phases:
\begin{itemize}
\item \textbf{Prepare:} proposer sends proposal number $N$.
\item \textbf{Promise:} acceptors promise not to accept lower numbers.
\item \textbf{Accept:} proposer asks acceptors to accept value.
\item \textbf{Learn:} majority acceptance becomes final value.
\end{itemize}

\begin{center}
\framebox[0.95\linewidth][c]{
\begin{minipage}{0.90\linewidth}
\vspace{0.1cm}
\centerline{\textbf{Paxos Flow (Simple)}}
\centerline{Prepare $\rightarrow$ Promise $\rightarrow$ Accept $\rightarrow$ Learn}
\vspace{0.1cm}
\end{minipage}}
\end{center}

\subsection{Proof of Work (PoW) and Mining}
PoW is used by Bitcoin to select who adds the next block.

Main idea:
\begin{itemize}
\item miners collect valid transactions,
\item build block candidate,
\item change nonce and hash repeatedly,
\item if hash is below target, block is valid.
\end{itemize}

Why useful:
\begin{itemize}
\item creating fake history becomes expensive,
\item verifying a valid block remains easy,
\item network gets security from real-world computation cost.
\end{itemize}

\begin{center}
\framebox[0.95\linewidth][c]{
\begin{minipage}{0.90\linewidth}
\vspace{0.1cm}
\centerline{\textbf{Bitcoin Mining Flowchart}}
\centerline{Pick TX from mempool}
\centerline{$\Downarrow$}
\centerline{Verify TX}
\centerline{$\Downarrow$}
\centerline{Build block + nonce}
\centerline{$\Downarrow$}
\centerline{Hash and compare with target}
\centerline{No: retry \hspace{0.6cm} Yes: broadcast block}
\vspace{0.1cm}
\end{minipage}}
\end{center}

\subsection{Proof of Stake (PoS)}
PoS chooses validators using stake (locked coins), not heavy computation.

\begin{itemize}
\item users lock coins to join as validators,
\item protocol selects one validator to propose block,
\item others verify and attest,
\item valid behavior gets rewards,
\item malicious behavior gets slashing penalty.
\end{itemize}

Quick comparison:
\begin{itemize}
\item \textbf{PoW:} leader by computation, high energy use, mature security history.
\item \textbf{PoS:} leader by stake, low energy use, faster and more efficient in many designs.
\end{itemize}

\subsection{Bitcoin Transaction and UTXO (Most Important)}
Bitcoin does not use account balance in the same way as a bank app.
It uses \textbf{UTXO} (Unspent Transaction Output).

Think like this:
\begin{itemize}
\item each UTXO is a spendable digital coin piece,
\item transaction inputs consume old UTXOs,
\item transaction outputs create new UTXOs,
\item one UTXO cannot be spent twice.
\end{itemize}

\begin{center}
\framebox[0.95\linewidth][c]{
\begin{minipage}{0.90\linewidth}
\vspace{0.1cm}
\centerline{\textbf{Bitcoin Transaction Structure}}
\centerline{Input(s): previous UTXO references + signature}
\centerline{$\Downarrow$}
\centerline{Transaction core data}
\centerline{$\Downarrow$}
\centerline{Output(s): recipient + change}
\vspace{0.1cm}
\end{minipage}}
\end{center}

\subsection{How Transaction Validation Works}
Nodes verify a transaction before accepting it:
\begin{itemize}
\item signatures must be valid,
\item referenced UTXOs must exist,
\item referenced UTXOs must be unspent,
\item output amount must not exceed input amount,
\item scripts/conditions must pass.
\end{itemize}

\subsection{Wallets and Types}
A wallet stores \textbf{keys}, not actual coins.
Coins remain on blockchain ledger.

\begin{itemize}
\item \textbf{Hot wallets:} online wallets (mobile/web/desktop), easy to use, higher attack surface.
\item \textbf{Cold wallets:} offline wallets (hardware/paper), safer for long-term storage.
\item \textbf{Custodial:} third party holds keys.
\item \textbf{Non-custodial:} user holds keys and full responsibility.
\end{itemize}

\examtip{Good practical strategy: hot wallet for small daily amount, cold wallet for larger long-term amount.}

\section{Solved Question Papers}

\subsection{Part A: 3-Mark Questions}

\subsubsection*{Q1. Describe the two main categories of consensus mechanisms.}
\begin{itemize}
\item Consensus is the rule that makes all nodes agree on one ledger state.
\item CFT handles crashes/timeouts and assumes nodes do not lie.
\item BFT handles malicious or conflicting node behavior.
\item Paxos is a CFT example; PoW and PoS are blockchain BFT-style approaches.
\end{itemize}

\subsubsection*{Q2. How are transactions validated in a Bitcoin network?}
\begin{itemize}
\item Node checks transaction structure and protocol rules.
\item Signature must prove sender authorization.
\item Inputs must reference valid unspent UTXOs.
\item Input value must be at least output value.
\item If valid, transaction is relayed and added to mempool.
\end{itemize}

\subsubsection*{Q3. Explain the working of PoW consensus mechanism.}
\begin{itemize}
\item Miner gathers transactions and forms a candidate block.
\item Miner keeps changing nonce and hashing block header.
\item Target is to get hash below difficulty limit.
\item First successful miner broadcasts block.
\item Network verifies quickly and adds block.
\item Miner receives block reward and fees.
\end{itemize}

\subsubsection*{Q4. What are the different steps involved in a Bitcoin transaction?}
\begin{itemize}
\item User enters receiver address and amount.
\item Wallet selects UTXOs and prepares outputs.
\item Sender signs with private key.
\item Transaction is broadcast to network.
\item Nodes validate and miners include it in a block.
\item Confirmations make payment safer to accept.
\end{itemize}

\subsubsection*{Q5. What is mining in blockchain?}
\begin{itemize}
\item Mining means validating transactions and creating blocks in PoW.
\item Miners solve a hash puzzle by nonce trial-and-error.
\item Winning miner publishes block to network.
\item Other nodes verify and append the block.
\item Mining gives security and transaction ordering.
\item Rewards and fees motivate miners to participate.
\end{itemize}

\subsubsection*{Q6. If network has 6 Byzantine nodes, what minimum nodes are required for PBFT?}
\begin{itemize}
\item PBFT uses formula $n \ge 3f+1$.
\item Here $f=6$.
\item So $n \ge 3(6)+1=19$.
\item Minimum required nodes are 19.
\end{itemize}

\subsubsection*{Q7. What is Byzantine generals problem and why relevant to blockchain?}
\begin{itemize}
\item It models agreement in presence of dishonest participants.
\item Honest nodes must still reach one common decision.
\item Blockchain faces this exact challenge in open networks.
\item Consensus protocols are designed to solve it.
\item Without this, ledger can split into conflicting histories.
\end{itemize}

\subsubsection*{Q8. What is a Sybil attack, and how does Bitcoin prevent it?}
\begin{itemize}
\item Sybil attack means one actor creates many fake nodes.
\item If voting is identity-based, attacker gains unfair control.
\item Bitcoin ties influence to hash power, not node count.
\item So fake identities alone do not help the attacker.
\item PoW makes control expensive due to hardware and energy cost.
\end{itemize}

\subsection{Part B: 14-Mark Questions}

\subsubsection*{Q1. Explain need for consensus algorithms. What are CFT algorithms? Compare PoW and PoS.}
\begin{itemize}
\item Consensus is needed because blockchain has no central authority.
\item It keeps all nodes on one valid ledger history.
\item It prevents double spending and conflicting states.
\item CFT handles crash failures, not malicious behavior.
\item Paxos is a well-known CFT protocol.
\item PoW selects leader by computational work.
\item PoS selects leader by stake and validator rules.
\item PoW is energy-heavy but proven.
\item PoS is energy-efficient and generally faster.
\item Both add economic cost to attack the network.
\end{itemize}

\subsubsection*{Q2. Explain Paxos protocol and give an implementation-style example.}
\begin{itemize}
\item Paxos solves consensus when nodes may crash or delay.
\item It defines proposer, acceptor, and learner roles.
\item Proposer starts with a numbered prepare message.
\item Acceptors promise to ignore lower proposal numbers.
\item Proposer sends accept request after majority promises.
\item Majority acceptance finalizes the value.
\item Learners record that final value.
\item In 5 nodes, majority is 3, so 2 crashes can be tolerated.
\item Paxos is safe for crash faults, not Byzantine attacks.
\item It is widely used in distributed backend systems.
\end{itemize}

\subsubsection*{Q3. Explain working of PoS with neat sketch. Compare PoW and PoS.}
\begin{itemize}
\item PoS works by locking coins as validator stake.
\item Protocol selects validator to propose a block.
\item Other validators verify and attest the block.
\item Enough attestations make the block final.
\item Honest validators earn rewards.
\item Dishonest validators can be slashed.
\item PoS uses less energy than PoW.
\item PoW relies on compute race and electricity.
\item PoW attack needs hash-power majority.
\item PoS attack needs very large stake with slashing risk.
\item[]
\vspace{0.25em}
\noindent\framebox[0.97\linewidth][c]{
\begin{minipage}{0.93\linewidth}
\vspace{0.08cm}
\centerline{\textbf{PoS Sketch}}
\centerline{Stake Lock $\rightarrow$ Validator Selection $\rightarrow$ Block Proposal}
\centerline{$\Downarrow$}
\centerline{Peer Verification $\rightarrow$ Finality $\rightarrow$ Reward/Slashing}
\vspace{0.08cm}
\end{minipage}}
\end{itemize}

\subsubsection*{Q4. Explain Bitcoin payment from user perspective. What is a wallet? Explain types.}
\begin{itemize}
\item User enters receiver address and amount.
\item Wallet selects UTXOs and creates outputs.
\item Transaction is signed with private key.
\item Signed transaction is broadcast to network.
\item Nodes validate and place it in mempool.
\item Miner includes it in a block.
\item Confirmations make payment increasingly final.
\item Wallet is a key-management tool, not coin storage.
\item Hot wallets are convenient but online risk exists.
\item Cold wallets are safer for long-term holdings.
\end{itemize}

\subsubsection*{Q5. Explain transaction validation, mining steps, and role of miner in Bitcoin network.}
\begin{itemize}
\item Nodes validate signatures, scripts, and input references.
\item Inputs must be unspent UTXOs with enough value.
\item Valid transactions are stored in mempool.
\item Miner selects transactions and builds a candidate block.
\item Miner performs nonce hashing against difficulty target.
\item On success, block is broadcast and verified by peers.
\item Accepted block updates ledger and UTXO set.
\item Miners provide ordering, validation, and network security.
\item Rewards and fees incentivize participation.
\item This process keeps Bitcoin decentralized and reliable.
\end{itemize}

\subsubsection*{Q6. Explain tasks of Bitcoin miner and mining algorithm with flowchart.}
\begin{itemize}
\item Miner gathers transactions from mempool.
\item Miner validates and filters invalid entries.
\item Miner forms candidate block and block header.
\item Miner performs repeated hashing with nonce.
\item Hash below target means valid proof.
\item Miner broadcasts block to the network.
\item Peers verify and append if valid.
\item Miner receives subsidy and fees.
\item Miner then starts next round.
\item Continuous competition keeps chain secure.
\item[]
\vspace{0.25em}
\noindent\framebox[0.97\linewidth][c]{
\begin{minipage}{0.93\linewidth}
\vspace{0.08cm}
\centerline{\textbf{Mining Algorithm Flowchart}}
\centerline{Collect TX $\rightarrow$ Validate TX $\rightarrow$ Build Header}
\centerline{$\Downarrow$}
\centerline{Hash + Nonce $\rightarrow$ Target Check}
\centerline{No: Retry \hspace{0.5cm} Yes: Broadcast Block}
\vspace{0.08cm}
\end{minipage}}
\end{itemize}

\subsubsection*{Q7. Random node (not PoW) is chosen. If chosen node is Byzantine, analyze stealing, DoS, and double-spend claims.}
\begin{itemize}
\item Random single-leader choice is weaker than PoW for security.
\item Byzantine leader can delay/censor transactions, so DoS is possible.
\item Direct theft of others' coins still needs their private keys.
\item Without private keys, direct stealing is false.
\item Reordering/censoring can support double-spend attempts.
\item Double-spend success depends on confirmations and protocol rules.
\item So double-spend attempt is true, guaranteed success is not.
\item Final answer: stealing false, DoS true, double-spend attempt true.
\item This is why Bitcoin uses expensive work-based leader selection.
\end{itemize}

\subsubsection*{Q8. Explain Bitcoin transaction structure with diagram. Explain UTXO and use in verification.}
\begin{itemize}
\item Bitcoin transaction contains inputs, outputs, and metadata.
\item Inputs reference previous outputs and provide authorization data.
\item Outputs define receiver value and lock conditions.
\item Unspent outputs are called UTXOs.
\item Validation checks input references exist and are unspent.
\item Validation checks signatures/scripts and amount consistency.
\item Input minus output becomes fee.
\item After confirmation, spent inputs are removed from UTXO set.
\item New outputs become fresh UTXOs.
\item UTXO model makes double-spend detection straightforward.
\item[]
\vspace{0.25em}
\noindent\framebox[0.97\linewidth][c]{
\begin{minipage}{0.93\linewidth}
\vspace{0.08cm}
\centerline{\textbf{Bitcoin Transaction Diagram}}
\centerline{Input 1: Prev TXID + Index + Unlock Script}
\centerline{Input 2: Prev TXID + Index + Unlock Script}
\centerline{$\Downarrow$}
\centerline{Transaction Data Core}
\centerline{$\Downarrow$}
\centerline{Output 1: Recipient Value + Lock Script}
\centerline{Output 2: Change Value + Lock Script}
\vspace{0.08cm}
\end{minipage}}
\end{itemize}

\end{document}
